\documentclass[11pt]{article}
\usepackage{amsmath,amssymb,amsthm,mathtools}
\usepackage{xcolor}
\usepackage[margin=1in]{geometry}
\usepackage{hyperref}
\hypersetup{colorlinks=true,linkcolor=blue!50!black,urlcolor=blue!50!black,citecolor=blue!50!black}

\newtheorem{theorem}{Theorem}
\newtheorem{lemma}{Lemma}
\newtheorem{corollary}{Corollary}
\newtheorem{proposition}{Proposition}
\theoremstyle{definition}
\newtheorem{definition}{Definition}
\theoremstyle{remark}
\newtheorem{remark}{Remark}

\newcommand{\LP}{\mathcal{LP}}
\newcommand{\C}{\mathbb{C}}
\newcommand{\R}{\mathbb{R}}
\newcommand{\Xibf}{\Xi}

\title{\bf A flow-blindness no-go:\\ soft hypotheses cannot prove the Riemann Hypothesis}
\author{mathexplorer research note (Lemma SL2)}
\date{2026-08}

\begin{document}
\maketitle

\begin{abstract}
We record a precise, unconditional obstruction to a large class of proposed proofs
of the Riemann Hypothesis (RH): those that would deduce RH from ``soft''
(analytic, structural) hypotheses on the de Bruijn--Newman heat profile
$\Phi$ --- evenness, positivity, doubly-exponential decay, log-concavity, total
positivity. Multiplying $\Phi$ by the Gaussian $e^{-\varepsilon u^2}$ is exactly
the de Bruijn backward heat flow at time $t=-\varepsilon$; it keeps $\Phi$ inside
every one of those soft classes, yet by the Rodgers--Tao theorem
$\Lambda\ge 0$ the resulting function has a non-real zero for \emph{every}
$\varepsilon>0$. Hence no property of $\Phi$ that is stable under its own
$\varepsilon$-backward-heat deformations can imply that its transform has only
real zeros. Any proof of RH must use a property distinguishing $\Phi$ from these
deformations --- necessarily an arithmetic property (the Jacobi theta / Euler
product identity behind $\Phi$), which the heat flow destroys. This upgrades
Newman's ``the Riemann hypothesis, if true, is only barely so'' from an aphorism
to a formal no-go, and is the analytic twin of the frozen-rung$\Leftrightarrow$
Platt--Trudgian inertia result: it re-derives, from de Bruijn--Newman flow
rigidity rather than from Sylvester inertia, the conclusion that the only genuine
lever in the program is verified zero height.
\end{abstract}

\section{Setup}

\paragraph{The de Bruijn--Newman flow.} Let $\Phi\colon\R\to\R$ be the standard
super-exponentially decaying de Bruijn--Newman profile,
\[
  \Phi(u)=\sum_{n\ge1}\bigl(2\pi^2 n^4 e^{9u}-3\pi n^2 e^{5u}\bigr)\,
    e^{-\pi n^2 e^{4u}},
\]
which is even, strictly positive, and decays doubly exponentially. For $t\in\R$
the de Bruijn heat flow of $\Phi$ is
\begin{equation}\label{eq:Ht}
  H_t(z)\;=\;\int_0^\infty e^{t u^2}\,\Phi(u)\,\cos(zu)\,du,
\end{equation}
an entire function of order $1$, real and even on $\R$. The de Bruijn--Newman
constant $\Lambda$ is the threshold such that $H_t$ has \emph{only real zeros} if
and only if $t\ge\Lambda$ \cite{deBruijn1950,Newman1976}. The Riemann Hypothesis
is the statement $\Lambda\le 0$; since $H_0=\tfrac18\,\Xi(\,\cdot\,)$ with
$\Xi(\lambda)=8\int_0^\infty\Phi(u)\cos(2\lambda u)\,du = c\,\xi(\tfrac12+i\lambda)$
the completed Riemann $\xi$, RH is exactly ``$H_0$ has only real zeros.''

\paragraph{The Rodgers--Tao theorem.} Rodgers and Tao \cite{RodgersTao2020} proved
\begin{equation}\label{eq:RT}
  \Lambda\;\ge\;0 .
\end{equation}
Equivalently: for every $t<0$, $H_t$ has at least one non-real zero. (Newman
\cite{Newman1976} had shown $\Lambda>-\infty$ and conjectured \eqref{eq:RT}.)

\paragraph{The soft class.} Write $\mathcal S$ for the class of admissible profiles
$\Psi\colon\R\to\R_{\ge0}$ that are even, strictly positive, log-concave, and decay
faster than $e^{-c u^2}$ for every $c$ (doubly-exponential decay), together with any
of the usual structural refinements (total positivity / P\'olya-frequency of the
associated kernel, moment positivity, etc.). We call a hypothesis on the profile
\emph{soft} if it defines a subclass $\mathcal H\subseteq\mathcal S$; and RH-\emph{stated-softly}
is the assertion ``$\Psi\in\mathcal H \Rightarrow$ the transform of $\Psi$ has only
real zeros.'' Every hypothesis used in the de Branges / twin-ladder positivity
program (kernel positive-definiteness, log-concavity, Wronskian total positivity,
decay, functional symmetry) is soft in this sense.

\section{The backward flow stays soft but leaves the real-zero class}

\begin{lemma}[Backward heat flow preserves $\mathcal S$]\label{lem:soft}
For $\varepsilon>0$ set $\Phi_\varepsilon(u):=e^{-\varepsilon u^2}\Phi(u)$. Then
$\Phi_\varepsilon\in\mathcal S$, and its cosine transform equals the backward flow,
\[
  \Xi_\varepsilon(\lambda):=8\!\int_0^\infty\!\Phi_\varepsilon(u)\cos(2\lambda u)\,du
  \;=\;8\,H_{-\varepsilon}(2\lambda).
\]
\end{lemma}

\begin{proof}
The Gaussian $e^{-\varepsilon u^2}$ is even, strictly positive, and log-concave
($-\varepsilon u^2$ concave), so the pointwise product $\Phi_\varepsilon$ is even,
positive, and log-concave; its decay is at least that of $\Phi$. Total positivity is
preserved because the Gaussian is a P\'olya frequency function and the P\'olya
class is closed under multiplication. Thus $\Phi_\varepsilon\in\mathcal S$. The
transform identity is immediate from \eqref{eq:Ht} with $t=-\varepsilon$ and the
substitution $z=2\lambda$: $e^{-\varepsilon u^2}\Phi(u)$ under the cosine transform
is $H_{-\varepsilon}$.
\end{proof}

\begin{theorem}[Flow-blindness]\label{thm:main}
For every $\varepsilon>0$ the profile $\Phi_\varepsilon\in\mathcal S$ satisfies every
soft hypothesis that $\Phi$ does, yet $\Xi_\varepsilon$ has a non-real zero. In
particular $\Phi_\varepsilon$ is a counterexample to RH-stated-softly for the class
$\mathcal H=\mathcal S$, and for every soft subclass containing the family
$\{\Phi_\varepsilon:0<\varepsilon<\varepsilon_0\}$.
\end{theorem}

\begin{proof}
By Lemma~\ref{lem:soft}, $\Phi_\varepsilon\in\mathcal S$ and
$\Xi_\varepsilon=8\,H_{-\varepsilon}(2\,\cdot\,)$. Since $-\varepsilon<0\le\Lambda$
by \eqref{eq:RT}, $-\varepsilon<\Lambda$, so $H_{-\varepsilon}$ --- hence
$\Xi_\varepsilon$ --- has a non-real zero. The soft hypotheses satisfied by $\Phi$
are preserved by the multiplier $e^{-\varepsilon u^2}$ (Lemma~\ref{lem:soft}), so
$\Phi_\varepsilon$ satisfies them too while failing ``all zeros real.''
\end{proof}

\begin{corollary}[No soft proof of RH]\label{cor:nogo}
Let $\mathcal H\subseteq\mathcal S$ be any soft class that is invariant under the
backward-heat family, i.e.\ $\Psi\in\mathcal H\Rightarrow e^{-\varepsilon u^2}\Psi\in\mathcal H$
for all small $\varepsilon>0$. Then RH-stated-softly for $\mathcal H$ is
\emph{false}. Consequently no proof of RH whose hypotheses on the profile are
flow-invariant can be correct: a correct proof must invoke a property of $\Phi$ not
shared by its backward-heat deformations.
\end{corollary}

\begin{proof}
Immediate from Theorem~\ref{thm:main}: $\Phi_\varepsilon\in\mathcal H$ for small
$\varepsilon$ by invariance, and $\Xi_\varepsilon$ has a non-real zero, so the
implication defining RH-stated-softly fails on $\Phi_\varepsilon$.
\end{proof}

\begin{remark}[What the distinguishing property must be]
The property of $\Phi$ that its deformations $\Phi_\varepsilon$ lack is precisely the
one destroyed by the heat flow: the exact arithmetic identity
$\Xi=c\,\xi(\tfrac12+i\cdot)$ tying $\Phi$ to the Jacobi theta function and the Euler
product. Multiplication by $e^{-\varepsilon u^2}$ breaks the functional
equation's arithmetic origin while preserving all analytic softness. Thus any RH
proof through this apparatus must, at some step, use prime-side input (e.g.\ Weil
explicit-formula positivity), the classical device that converts prime information
into zero-location information. No purely soft/estimate mechanism suffices.
\end{remark}

\section{Consequence for the twin-ladder / de Branges ``T3'' strip ladder}

The program's sole global RH frontier (``T3'') asks, for every shift $a>0$, that a
finite Hermitian kernel $D_{a,R}+K_a^{(S)}+K_{a,0}$ be positive semidefinite. With
the repo transform pair
$A_a(z)=8\!\int_0^\infty\!\Phi\cosh(2au)\cos(2zu)\,du$,
$B_a(z)=8\!\int_0^\infty\!\Phi\sinh(2au)\sin(2zu)\,du$, one has the exact identity
(machine-verified to $10^{-30}$ in the trigonometric core)
\begin{equation}\label{eq:dagger}
  E_a(z):=A_a(z)-iB_a(z)=\Xi(z+ia),\tag{$\dagger$}
\end{equation}
so $E_a$ is the Hermite--Biehler function of the $a$-shifted $\Xi$ and the kernel is
its de Branges reproducing kernel \cite{deBranges1968}. Positivity of that kernel is
equivalent to $E_a\in$ Hermite--Biehler, i.e.\ to ``$\Xi$ has no zero with
$\Im>a$,'' i.e.\ to ``$\zeta$ has no zero with $\Re s>\tfrac12+a$.'' Thus T3 is the
zero-free-strip statement in operator language, and ``T3$(a)$ for all $a>0$'' is RH.

\begin{corollary}
The T3 strip ladder cannot be closed by any soft/estimate mechanism. Its hypotheses
(kernel positivity, Wronskian total positivity, decay, symmetry) are flow-invariant,
so by Corollary~\ref{cor:nogo} they hold for $\Phi_\varepsilon$, whose $E_a^{(\varepsilon)}$
fails Hermite--Biehler for all $a$ below some $b(\varepsilon)>0$. In particular no
cancellation-sensitive Volterra estimate uniform over flow-stable profiles can force
the strip: the family $\{\Phi_\varepsilon\}$ is an arbitrarily small deformation of
$\Phi$ that the estimate would also have to control, but does not.
\end{corollary}

\section{Relation to the rest of the program}

\paragraph{Newman, made formal.} Newman's dictum that RH ``if true is only barely
so'' \cite{Newman1976} is exactly \eqref{eq:RT}: the flow reaches all-real-zeros no
earlier than $t=0$. Theorem~\ref{thm:main} turns this boundary-case rigidity into a
proof-theoretic statement: the RH-true profile sits on the exact boundary of the
soft class's real-zero region, so no flow-stable (soft) condition can separate it
from the non-RH side.

\paragraph{Twin of frozen-rung$\Leftrightarrow$Platt--Trudgian.} The frozen-rung
result shows, via Sylvester inertia / Krein--Langer, that freezing the twin-ladder
rungs reproduces Platt--Trudgian in operator language --- the route can never beat
verified zero height. Theorem~\ref{thm:main} reaches the same ceiling from the
opposite side: the analytic surplus a soft argument would need is precisely a
zero-free strip, and no soft hypothesis can supply it. Both say: \emph{the only
genuine lever is verified height}, now derived independently from inertia and from
de Bruijn--Newman flow rigidity.

\paragraph{Scope and honesty.} Theorem~\ref{thm:main} and Corollary~\ref{cor:nogo}
are unconditional (they rest only on \cite{RodgersTao2020} and elementary closure
properties). They do \emph{not} say RH is false or unprovable; they say a specific,
large family of \emph{soft} strategies --- everything the twin-ladder/de~Branges
positivity apparatus currently uses --- cannot work, and locate what a proof must
add (prime-side arithmetic input). The value is negative-but-sharp: it retires the
sub-threshold residual-grinding program and redirects effort to arithmetic
mechanisms.

\begin{thebibliography}{9}
\bibitem{deBruijn1950} N.~G.~de Bruijn, \emph{The roots of trigonometric integrals},
  Duke Math.\ J.\ \textbf{17} (1950), 197--226.
\bibitem{Newman1976} C.~M.~Newman, \emph{Fourier transforms with only real zeros},
  Proc.\ Amer.\ Math.\ Soc.\ \textbf{61} (1976), 245--251.
\bibitem{RodgersTao2020} B.~Rodgers and T.~Tao, \emph{The de Bruijn--Newman
  constant is non-negative}, Forum Math.\ Pi \textbf{8} (2020), e6.
\bibitem{deBranges1968} L.~de Branges, \emph{Hilbert spaces of entire functions},
  Prentice--Hall, 1968.
\end{thebibliography}

\end{document}
